\section{A2 - Object Detection và phân tích paper}

\subsection{Bài toán, dataset và A2.1 protocol}
Chủ đề đã chọn là \textbf{Object Detection}. Dataset đề xuất hiện tại là BDD100K detection vì có road scenes, nhiều class và object ở nhiều scale. Trước A2.1 cần chốt dataset/subset, class mapping, split, metric và paper citation.

\TODO{XÁC NHẬN DATASET A2, QUY MÔ SUBSET VÀ PAPER CUỐI CÙNG TRƯỚC 02/11/2026.}
\safeimage{a2/a2_eda.png}{EDA A2: số box theo lớp, số object/ảnh, phân bố diện tích/aspect ratio và sample annotations.}

\subsection{Dataset, DataLoader và augmentation}
Mỗi sample trả về ảnh và target gồm boxes, labels, area, image id và iscrowd. Ảnh được resize về một canvas cố định cho detector; bounding boxes được scale tương ứng. Horizontal flip hoặc augmentation hình học chỉ được dùng khi code cập nhật box chính xác. Validation/test không dùng augmentation ngẫu nhiên.

\subsection{Paper chính: Feature Pyramid Networks for Object Detection}
Paper chính dự kiến là Lin et al., \emph{Feature Pyramid Networks for Object Detection} \cite{lin2017fpn}; Faster R-CNN \cite{ren2015faster} là detector nền. Phần paper phải viết lại bằng lời của nhóm, không chép abstract, và gồm đủ method, experiments, limitations và liên hệ với pipeline.

\subsubsection{Ý tưởng, kiến trúc và giả thiết}
FPN dùng top-down pathway và lateral connections để tạo pyramid feature có semantic mạnh ở nhiều resolution. Backbone tạo $C_2,C_3,C_4,C_5$; mỗi level được chiếu bằng convolution $1\times1$, cộng với feature cấp cao được upsample, rồi convolution $3\times3$ tạo $P_2,P_3,P_4,P_5$.
\safeimage{a2/a2_fpn_architecture.png}{Top-down pathway và lateral connections của custom FPN.}

\subsubsection{Thí nghiệm của paper}
\TODO{ĐỌC PAPER CHÍNH THỨC VÀ ĐIỀN: DATASET, BACKBONE, DETECTOR, METRIC, ABLATION, KẾT LUẬN; CITE TỪNG CLAIM QUAN TRỌNG.}

\subsubsection{Hạn chế của paper và khả năng tái hiện}
Báo cáo phải phân biệt rõ \emph{full benchmark reproduction} với \emph{controlled/partial reproduction}. Nếu nhóm dùng BDD100K/subset, schedule khác và pretrained backbone thì không được tuyên bố tái hiện đúng số liệu COCO của paper; mục tiêu là kiểm tra lại contribution FPN trong protocol có kiểm soát.

\subsection{Pipeline Faster R-CNN tự ghép theo từng thành phần}
Repository không gọi \code{torchvision.models.detection.FasterRCNN} và không gọi \code{fasterrcnn\_resnet50\_fpn}. Pipeline được viết thành các module độc lập:
\begin{enumerate}
\item fixed-size detection transform và box scaling;
\item ResNet stages và single-scale/FPN backbone;
\item custom anchor generator;
\item custom RPN head, anchor matching, sampling, box regression và proposal generation;
\item RoIAlign là operator cấp thấp từ torchvision;
\item custom RoI sampling, two-FC box head và class-specific box predictor;
\item custom training loss assembly và inference post-processing;
\item NMS là operator tối ưu cấp thấp từ torchvision.
\end{enumerate}
\safeimage{a2/a2_pipeline.png}{Pipeline Faster R-CNN/FPN được hiện thực trong repository.}

\subsection{Ma trận thí nghiệm A2}
\begin{table}[H]\centering
\caption{Bốn biến thể A2.}
\begin{tabularx}{\textwidth}{@{}llXX@{}}
\toprule
ID & Detector & Biến thay đổi & Câu hỏi \\
\midrule
D1 & Faster R-CNN single-scale & Chỉ C5 projection, không FPN & Baseline detection \\
D2 & Faster R-CNN + FPN & Custom P2--P5 & FPN cải thiện multi-scale detection thế nào? \\
D3 & FPN + small anchors & Anchor scales nhỏ hơn & Có tăng AP của small objects không? \\
D4 & FPN + deeper fine-tune & Mở thêm ResNet stages & Accuracy/compute trade-off thế nào? \\
\bottomrule
\end{tabularx}
\end{table}

\subsection{Training protocol}
Backbone có thể khởi tạo ImageNet pretrained; detector heads/FPN được train trên dataset A2. Không bắt buộc train toàn bộ backbone từ random initialization. Để D1--D2 là controlled comparison, cần giữ giống nhau split, augmentation, image size, optimizer, schedule và backbone initialization; biến chính là FPN.

\begin{table}[H]\centering
\caption{Thiết lập huấn luyện A2.}
\begin{tabularx}{0.94\textwidth}{@{}lX@{}}
\toprule
Thành phần & Giá trị \\
\midrule
Train/val/subset & \TBD \\
Image size & \TBD \\
Backbone init / trainable stages & \TBD \\
Optimizer / LR / schedule & \TBD \\
RPN anchors / NMS thresholds & \TBD \\
Model selection & validation mAP \\
\bottomrule
\end{tabularx}
\end{table}

\subsection{Kết quả định lượng}
\begin{table}[H]\centering
\caption{Kết quả chính A2.}
\begin{tabular}{@{}lrrrrrr@{}}
\toprule
ID & mAP@0.5 & mAP@0.5:0.95 & AP$_S$ & AP$_M$ & AP$_L$ & Latency \\
\midrule
D1 & \metricpending & \metricpending & \metricpending & \metricpending & \metricpending & \metricpending \\
D2 & \metricpending & \metricpending & \metricpending & \metricpending & \metricpending & \metricpending \\
D3 & \metricpending & \metricpending & \metricpending & \metricpending & \metricpending & \metricpending \\
D4 & \metricpending & \metricpending & \metricpending & \metricpending & \metricpending & \metricpending \\
\bottomrule
\end{tabular}
\end{table}
\safeimage{a2/a2_pr_ap.png}{Precision-recall/AP comparison giữa D1--D4.}
\safeimage{a2/a2_detections.png}{Qualitative detections: đúng, false positive, false negative, small/occluded objects.}

\subsection{Failure-case analysis}
Phân nhóm lỗi theo small objects, occlusion, dense traffic, day/night, weather và class confusion. Các ví dụ phải được lấy từ output thật và đi kèm confidence/ground-truth để tránh cherry-picking.

\subsection{Đối chiếu paper với pipeline của nhóm}
\begin{table}[H]\centering
\caption{Khung đối chiếu paper FPN và thí nghiệm nhóm.}
\begin{tabularx}{\textwidth}{@{}lXX@{}}
\toprule
Yếu tố & Paper & Pipeline nhóm \\
\midrule
Dataset & \TODO{ĐIỀN} & \ATwoDataset \\
Backbone/init & \TODO{ĐIỀN} & \TODO{ĐIỀN CONFIG THẬT} \\
Detector/FPN & \TODO{ĐIỀN} & Custom Faster R-CNN + custom FPN \\
Training schedule & \TODO{ĐIỀN} & \TODO{ĐIỀN} \\
Metric & \TODO{ĐIỀN} & mAP@0.5, mAP@0.5:0.95, AP by size \\
Kết luận & \TODO{ĐIỀN} & \TODO{ĐIỀN SAU D1--D4} \\
\bottomrule
\end{tabularx}
\end{table}

\subsection{Hạn chế và kết luận A2}
\TODO{THẢO LUẬN GIỚI HẠN GPU 12GB, SUBSET NẾU CÓ, PRETRAINED BACKBONE, CLASS IMBALANCE, SMALL OBJECTS, WEATHER/TIME-OF-DAY DOMAIN SHIFT, VÀ MỨC ĐỘ REPRODUCTION THỰC SỰ.}
